\subsection{Factorial Wick growth and the present implementation boundary}
\label{subsec:future_wick_scaling}

The present compiler treats every complete flavor-compatible Wick pairing as an
explicit symbolic contribution before PhysicsIR and the numerical contraction
graphs are constructed. This has the advantage of making the fermionic
semantics particularly transparent and provides a useful reference
implementation, but Eq.~\eqref{eq:number_wick_pairings} also exposes its main
scaling limitation. For repeated flavors the number of pairings grows as
\(\prod_f n_f!\). Thus a proton-like \(uud\) two-point function has two
same-flavor pairings, a \(pn\) system with three \(u\) and three \(d\) quarks has
\(3!3!=36\), and a \(ppn\) system contains \(5!4!=2880\) flavor-compatible
pairings before any numerical reuse is considered. Global common-subexpression
analysis can eliminate repeated numerical work after these terms have been
constructed, but it does not remove the factorial symbolic expansion itself.

This distinction is important for interpreting the current scope of SymCorr.
The explicit Wick engine is exact and useful as a correctness oracle, but it
should not be assumed to be the final representation for large multi-baryon
systems. A future representation should preserve the same flavor conservation,
field ordering, fermionic signs, and operator structures while avoiding the
need to materialize every permutation independently.

\subsection{Determinant compression and what it does not remove}
\label{subsec:future_determinants}

For fixed source and sink one-particle labels, the sum over identical-flavor
permutations has an exact determinant representation. If \(a_i\) and \(b_j\)
denote fixed sink and source labels for \(n_f\) quarks of flavor \(f\), then

\begin{equation}
\sum_{\sigma\in S_{n_f}}
\operatorname{sgn}(\sigma)
\prod_{i=1}^{n_f}
S_f(a_i,b_{\sigma(i)})
=
\det M_f,
\qquad
(M_f)_{ij}=S_f(a_i,b_j).
\label{eq:future_determinant_identity}
\end{equation}

This is the algebraic basis of determinant approaches to nuclear correlation
functions~\cite{PhysRevD.110.114505}. The determinant compresses the
\(n_f!\) fermionic permutations for a fixed set of scalar propagator entries.
It does not, by itself, eliminate sums over operator coefficients, spatial
coordinates, LapH modes, derivatives, or independent hadron momenta.

This point is especially important in a distillation calculation. The
perambulator \(\tau_f\) is an operator in LapH-mode and spin space. A baryon
spatial kernel such as Eq.~\eqref{eq:three_quark_spatial_kernel} couples three
LapH indices through a common spatial sum and the color tensor
\(\epsilon_{abc}\). In general it is therefore not a product of three
independent one-quark vectors. Consequently, taking a determinant of the
perambulator itself is not a replacement for the complete baryon or
multi-baryon contraction. Such a construction would compress the fermionic
permutation structure while leaving, or potentially expanding, the dense
LapH-mode summations that the tensor formulation is designed to evaluate
collectively.

The same observation limits simple low-rank or Slater-determinant
factorizations of a generic baryon elemental. No such factorization is assumed
by the current compiler, and the future optimization strategy should not rely
on a generic three-quark LapH tensor being separable into independent one-body
factors. Where an operator does possess an additional exact factorization, it
can be exploited as an optimization, but it should be proved from the operator
structure rather than inferred from the fact that the final state contains
identical fermions.

\subsection{Reconstructing the projected coordinate-space propagator}
\label{subsec:future_reconstructed_propagator}

There is, however, another exact ordering of the distillation contractions in
which the LapH sums are removed before the determinant is formed. Combining
the LapH eigenvectors with the perambulator gives the distillation-projected
coordinate-space propagator

\begin{equation}
S_f^{\mathrm{sm}}(t,t_0)
=
V(t)\,\tau_f(t,t_0)\,V^\dagger(t_0)
=
\Box(t)\,M_f^{-1}(t,t_0)\,\Box(t_0).
\label{eq:future_reconstructed_smeared_propagator}
\end{equation}

This operation does not undo distillation and does not reconstruct the full
unsmeared all-to-all propagator. It is the exact propagator inside the retained
projected subspace. Once the source and sink coordinate, color, and spin labels
are fixed, its matrix elements are ordinary scalars and may therefore enter the
determinant identity in Eq.~\eqref{eq:future_determinant_identity} directly.
Field-local transformations lead to corresponding variants such as
\((D V)\tau V^\dagger\), \(V\tau(DV)^\dagger\), or the analogous transported or
profiled constructions.

This representation can be attractive for compact local operators when only a
small set of source positions or endpoint transformations is required. The
same reconstructed propagator can then be reused by many local color--spin
operator structures, so increasing the final correlation-matrix dimension does
not necessarily require reconstructing the propagator again. In this regime a
determinant formulation and a native LapH formulation are best viewed as two
exact elimination orders of the same projected tensor network.

The limitation is the coordinate-space size of the reconstructed object. If
\(N_x\) sink sites and \(N_y\) source sites are retained, storage scales as
\(O(N_xN_y)\) times the color--spin matrix size. A fixed or small source set can
therefore be practical, whereas materializing all \(x,y\) pairs scales with the
square of the spatial volume and can erase much of the compression that makes
distillation useful. The same difficulty appears for scattering operators
with independently momentum-projected constituents: several independent
relative coordinates remain, and reconstructing coordinate-space propagation
before performing those sums can create a much larger problem than first
forming momentum-projected LapH elementals.

For this reason the reconstructed-propagator route should not be treated as a
universal replacement for native distillation. Compact local operators and
multi-hadron scattering operators expose different favorable contraction
orders even though both are exact descriptions of the same projected theory.

\subsection{Clustered antisymmetry and flavor-flow reduction}
\label{subsec:future_flow_compression}

A second exact direction retains the LapH representation but compresses the
fermionic permutation structure before all Wick terms are expanded. Suppose a
composite operator can be partitioned into spatial--color connected clusters
using the same incidence information already used to construct operator
vertices. For flavor \(f\), let \(r_a^{(f)}\) be the number of quarks in sink
cluster \(a\), and let \(c_b^{(f)}\) be the corresponding number of source
antiquarks in cluster \(b\). An inter-cluster flavor flow is described by a
non-negative integer matrix \(K^{(f)}\) satisfying

\begin{equation}
\sum_b K_{ab}^{(f)}=r_a^{(f)},
\qquad
\sum_a K_{ab}^{(f)}=c_b^{(f)}.
\label{eq:future_flavor_flow_constraints}
\end{equation}

Permutations that differ only by exchanges of identical-flavor fields inside
one source or sink cluster can then be grouped into the same exact flow class.
In group-theoretic language, the classes are double cosets of the complete
same-flavor permutation group by the local permutation subgroups. The local
antisymmetrization must act on the complete field slots: exchanging two
identical quarks exchanges their mode, spin, derivative or displacement role,
and all other field-local information together. A permutation of LapH axes
alone would not be the correct fermionic operation.

Exploratory exact algebraic tests motivate this direction. For the usual
nucleon flavor compositions, the explicit pairing counts \(2\), \(36\), and
\(2880\) for one nucleon, \(pn\), and \(ppn\), respectively, can be organized
into \(1\), \(4\), and \(77\) structural flavor-flow classes for the corresponding
cluster decomposition. These numbers describe an exact combinatorial
reorganization; they are not measured SymCorr performance results and the
compressed representation is not part of the present production compiler.
After such a reduction, the remaining flow graphs would still require ordinary
LapH tensor contraction and global reuse analysis. In particular, highly
exchanged three-cluster topologies can remain computationally expensive even
after the factorial permutation count has been reduced.

This suggests a useful future division of labor. Fermion algebra should remove
permutation redundancy that is exact at the symbolic level, while the existing
tensor optimizer should continue to decide contraction order, intermediate
materialization, and numerical reuse for the surviving LapH graphs. Such a
design would retain one generic compiler workflow rather than introducing
separate algorithms selected by hadron names.

\subsection{Representation choice and operator-basis reuse}
\label{subsec:future_representation_choice}

The preceding alternatives show that the nominal number of operators in a
correlation matrix is not by itself a sufficient predictor of cost. A basis of
many final operators can be inexpensive if those operators are linear
combinations of a small set of repeated spatial, derivative, momentum, and spin
primitives. Conversely, a smaller basis can be expensive if each operator
requires genuinely different spatial geometry or endpoint transformations. A
useful future cost model should therefore distinguish the number of final
operators from the number of unique expensive primitives and should account for
reuse before estimating matrix-wide cost.

Table~\ref{tab:future_representation_tradeoffs} summarizes the qualitative
tradeoffs identified above. The entries are intended as design guidance, not
as benchmark claims.

\begin{table}[t]
\centering
\caption{Qualitative comparison of exact contraction representations under
consideration for future SymCorr development. ``Current'' denotes the
explicit-Wick/native-LapH route implemented by the present compiler.}
\label{tab:future_representation_tradeoffs}
\begin{tabular}{p{0.23\linewidth}p{0.29\linewidth}p{0.35\linewidth}}
\toprule
Representation & Main strength & Principal limitation \\
\midrule
Explicit Wick + native LapH & Transparent reference semantics; spatial sums are compressed into reusable LapH objects & Factorial Wick-term generation for many identical flavors \\
Reconstructed \(V\tau V^\dagger\) + determinants & Removes LapH sums before determinant evaluation; potentially strong reuse for compact local operators & Coordinate-space storage and work grow rapidly with many source positions, relative coordinates, or endpoint transformations \\
Flow-compressed native LapH & Removes exact local permutation redundancy while retaining momentum-projected LapH geometry & Remaining flow tensor networks can still have high-rank mode-space contractions \\
\bottomrule
\end{tabular}
\end{table}

A long-term compiler could therefore preserve several mathematically equivalent
representations long enough to estimate their cost and reject infeasible ones.
The decision should depend on quantities such as \(N_v\), spatial volume,
number of independent position sums, number of source positions, number of
endpoint transformations, number of unique operator primitives, reuse across
the correlation matrix, and available memory. It should not be encoded as a
rule such as ``use determinants for nuclei'' or ``use LapH tensors for
baryons.'' Hardware-dependent choices should remain below this representation
boundary, in the execution-planning layers already responsible for memory and
device feasibility.

\subsection{Claims deliberately not made in the present work}
\label{subsec:future_nonclaims}

The considerations above also define several limits on what can presently be
claimed. The current SymCorr implementation does not contain the flavor-flow
compiler described above, does not automatically reconstruct
\(V\tau V^\dagger\) as an alternative determinant representation, and does not
select between several algebraically equivalent physics representations using a
cost model. The determinant identity does not imply that a determinant of a
perambulator alone is a complete baryon correlator, and no generic low-rank
factorization of dense baryon LapH elementals is assumed. Likewise, the
preliminary reduction in the number of flow classes should not be interpreted
as an equivalent wall-clock speedup before the surviving tensor contractions
are benchmarked in the full runtime.

No approximate pruning of Wick terms is currently part of the exact compiler.
Possible future uses of randomized polynomial identity testing, determinant
reuse, approximate low-rank structure, or magnitude-based truncation require
separate validation and should remain distinct from exact algebraic
canonicalization. The immediate development priority is therefore narrower:
retain the explicit Wick path as a reference, introduce any compressed fermion
representation behind an independently testable compiler boundary, prove exact
equivalence on small systems, and only then compare the competing numerical
representations with controlled CPU/GPU benchmarks.
